\titulocentrado{Resumen}

La gestión administrativa de los salones de belleza de pequeña escala suele apoyarse en hojas de cálculo elaboradas manualmente, lo que favorece errores de cálculo, inconsistencias de datos y una liquidación de sueldos lenta. El presente trabajo plantea el diseño y la validación de un sistema web móvil, denominado \textit{Lush Beauty}, orientado a la gestión de caja y a la liquidación semanal de comisiones del salón homónimo, cuyo diagnóstico se basa en una planilla de pagos de ejemplo, representativa de la información que recibe el salón (agosto de 2026) (91 servicios, tres especialistas). La investigación es de enfoque cuantitativo con apoyo cualitativo, alcance descriptivo-explicativo y diseño preexperimental de preprueba y posprueba. Se incluyen el planteamiento del problema, los objetivos, la justificación, el marco teórico, un benchmarking de sistemas comerciales similares, las hipótesis con su operacionalización, el marco metodológico y el cronograma de actividades. Se espera que el sistema reduzca el tiempo de liquidación semanal y la tasa de errores de cálculo, además de alcanzar una usabilidad percibida aceptable por parte de las usuarias. Como ampliación orientada a la toma de decisiones, se propone evolucionar el sistema hacia un Sistema de Soporte para la Toma de Decisiones, denominado Lush Beauty Smart DSS, que conserva el núcleo transaccional de caja y nómina y agrega análisis de contribución por servicio, simulación de escenarios What-If, comparación multicriterio con análisis de sensibilidad, recomendaciones explicables y seguimiento de decisiones, además de registro sin conexión, importación asistida de Excel, un asistente financiero (Lush Insight), alertas (Lush Guard) y automatizaciones con n8n. La decisión final siempre corresponde a la administradora: la inteligencia artificial no autoriza pagos ni cambia por sí sola las políticas económicas, y los datos ilustrativos o los escenarios hipotéticos no se presentan como resultados medidos.

\noindent\textbf{Palabras clave:} sistema de información, gestión de salones de belleza, liquidación de comisiones, aplicación web progresiva, benchmarking, sistemas de soporte para decisiones, simulación What-If, análisis multicriterio, inteligencia artificial explicable.

\titulocentrado{Abstract}

\selectlanguage{english}
Lush Beauty Smart DSS extends a mobile-first cash and payroll management system for a small beauty salon into a decision support platform. It integrates service-level contribution analysis, What-If simulation, transparent weighted decision comparison with sensitivity testing, an explainable recommendation engine, and a decision outcome log. The proposal retains offline-first operations, reviewed spreadsheet imports, finance-focused AI assistance, rule-based anomaly alerts and n8n automations. Business decisions remain under human control. The illustrative dataset in the source report is not treated as observed business impact. This document defines architecture, assumptions, a worked example, testable requirements, evaluation metrics and a phased academic implementation plan.

\noindent\textbf{Keywords:} decision support systems; multi-criteria analysis; What-If simulation; contribution analysis; explainable artificial intelligence; cash management; commissions; progressive web app; n8n.
\selectlanguage{spanish}
